\chapter{Recurrence, Nonequilibrium Order, and Ising Spins}

\lectureoverview{A gas released from one half of a room rushes toward
equilibrium, yet reversible mechanics says that sufficiently long evolution
must also produce returns to the exceptional state.  We estimate that return
time from phase-space volume, separate recurrence from the unexplained
low-entropy beginning of the universe, and use the resulting puzzle to discuss
Boltzmann brains, coherent history, and life maintained by an energy flow.
We then turn to magnetic order.  Independent two-state spins give an exactly
soluble warm-up: combinatorics determines the partition function and the
magnetization is a hyperbolic tangent.  Neighbor interactions finally restore
an up--down symmetry and lead to the one-dimensional Ising model.}

\section{A Roomful of Air That Eventually Returns}

Begin with all the air in the left half of a room.  One can prepare this state
with a partition: evacuate the right half, place the air on the left, and then
remove the wall.  The gas spreads rapidly through the room and approaches
thermal equilibrium.  Entropy increases.  Nothing surprising has happened.

Now leave the room sealed and wait.  A million consecutive heads is an
extraordinarily unlikely coin-toss record, but an indefinitely repeated
experiment eventually produces extraordinarily unlikely records.  The same
logic applies to the gas.  At some remote time all of its molecules will again
be found in the left half, or in some comparably special region.  Such a return
is a Poincare recurrence.

The precise theorem is slightly more careful than the picture.  For a
finite-measure region of phase space whose dynamics preserves phase-space
volume, almost every initial point returns arbitrarily close to its starting
point.  A macroscopic statement such as \emph{all molecules are in the left
half} is easier: it asks only that the trajectory re-enter a finite
coarse-grained region, not that every molecular coordinate return exactly.
The practical question is how long that takes.

\subsection{A high-dimensional phase point}

For \(N\) classical particles, phase space has \(6N\) coordinates,
\begin{equation}
 \Gamma=(\bm x_1,\ldots,\bm x_N;\bm p_1,\ldots,\bm p_N).
 \label{eq:sm8-phase-point}
\end{equation}
The positions are bounded by the room.  At fixed total energy the momenta are
bounded as well, because one particle cannot carry arbitrarily large momentum
without carrying arbitrarily large energy.  The allowed energy shell
therefore has finite measure under the assumptions of the recurrence
argument.

The phase point moves deterministically, but in a many-body gas its motion is
chaotic.  Nearby trajectories separate quickly, as they do in a sequence of
billiard-ball collisions.  After coarse-graining, a long trajectory appears
to wander throughout the accessible region.

\begin{figure}[t]
\centering
\includegraphics[width=0.90\textwidth]{lecture_08_figure_02.png}
\caption{A schematic phase-space region and a complicated trajectory
wandering through it.  The drawing is deliberately low-dimensional; the real
gas lives in a \(6N\)-dimensional space.}
\lectureframe{00:06:41}
\end{figure}

The low-dimensional sketch immediately tempts us into a mistake.  If the
special state occupied half of the drawn rectangle, one might say that the
gas should spend half its time there.  But \emph{every molecule lies in the
left half} constrains every position coordinate, not just one coordinate of
the phase point.

Take only two particles moving in one spatial dimension, with coordinates
\(x_1\) and \(x_2\).  Their configuration space is a square.  Requiring
\(x_1<L/2\) selects half the square; requiring both
\(x_1<L/2\) and \(x_2<L/2\) selects one quarter:
\begin{equation}
 P_2=\frac{1}{2}\frac{1}{2}=\frac{1}{4}.
 \label{eq:sm8-two-particle-quarter}
\end{equation}

\begin{figure}[t]
\centering
\includegraphics[width=0.90\textwidth]{lecture_08_figure_07.png}
\caption{For two one-dimensional particles, the smaller square represents
the configurations in which both coordinates lie in the chosen half of the
room.}
\lectureframe{00:07:08}
\end{figure}

Three particles give \(1/2^3\), and \(N\) particles give
\begin{equation}
 P_{\mathrm{half}}(N)
 =\left(\frac{1}{2}\right)^N.
 \label{eq:sm8-half-probability}
\end{equation}
The exponent is the entire point.  A factor of one half is ordinary; one half
raised to a macroscopic particle number is beyond ordinary intuition.

\begin{classroomqa}
\audiencequestion{Should the answer be \(1/3^N\), since the molecules also
have other spatial directions?}
\lecturerresponse{It depends on the condition being imposed.  If the only
question is whether every particle lies to the left of the dividing wall,
then each particle satisfies one binary condition and the answer is
\(2^{-N}\).  Constraining another independent fraction of its available
position space would supply another factor.
\lecturetimestamp{00:07:47--00:08:05}}
\end{classroomqa}

\section{Phase Volume, Entropy, and the Recurrence Scale}

The same estimate can be written without referring to halves.  Let the room
have ordinary spatial volume \(V\), and let \(v\) be a smaller region in which
we demand that every molecule be found.  Ignoring momentum factors that are
common to the two macrostates, the available configuration-space volumes are
\begin{equation}
 \Omega_{\mathrm{room}}\propto V^N,
 \qquad
 \Omega_{\mathrm{small}}\propto v^N.
 \label{eq:sm8-config-volumes}
\end{equation}
The long-time fraction spent in the small macrostate is therefore
\begin{equation}
 P(v)
 =\frac{\Omega_{\mathrm{small}}}{\Omega_{\mathrm{room}}}
 =\left(\frac{v}{V}\right)^N.
 \label{eq:sm8-volume-ratio}
\end{equation}

\begin{figure}[t]
\centering
\includegraphics[width=0.90\textwidth]{lecture_08_figure_08.png}
\caption{The board's general phase-volume picture: a small region \(v^N\)
inside an accessible region \(V^N\), with the ratio \(v^N/V^N\).}
\lectureframe{00:10:44}
\end{figure}

In units with \(k_B=1\), entropy is the logarithm of accessible phase volume.
If \(S_A=\log\Omega_A\) is the entropy of the constrained macrostate and
\(S_{\rm eq}=\log\Omega_{\rm eq}\) that of equilibrium, then
\begin{equation}
 \begin{aligned}
 P(A)
 &=\frac{\Omega_A}{\Omega_{\rm eq}}\\
 &=\exp\!\left(S_A-S_{\rm eq}\right)\\
 &=e^{-\Delta S},\\
 \Delta S&=S_{\rm eq}-S_A.
 \end{aligned}
 \label{eq:sm8-entropy-probability}
\end{equation}
For the half-room state, \(\Delta S=N\log 2\), and
\eqref{eq:sm8-entropy-probability} reduces exactly to \(2^{-N}\).
The shorthand \(P\sim e^{-S_{\rm eq}}\) is appropriate only when the entropy
of the special region is being treated as a fixed reference; the physically
meaningful exponent is the entropy deficit \(\Delta S\).

\begin{classroomqa}
\audiencequestion{Is this the volume of configuration space rather than the
whole phase space?}
\lecturerresponse{For this comparison, yes.  We are imposing a positional
condition while leaving the momentum distribution at the same temperature,
so the common momentum-space factor cancels.  A fully general recurrence
statement uses phase space; this particular probability can be computed from
the configuration-space ratio.
\lecturetimestamp{00:11:21--00:11:40}}
\end{classroomqa}

Suppose the system loses memory of its detailed position on a microscopic
correlation or mixing time \(\tau_{\rm mix}\).  A rare region of probability
\(P(A)\) is then encountered on the scale
\begin{equation}
 \tau_{\rm rec}
 \sim \frac{\tau_{\rm mix}}{P(A)}
 \sim \tau_{\rm mix}e^{\Delta S}.
 \label{eq:sm8-recurrence-time-general}
\end{equation}
For the half-room condition,
\begin{equation}
 \tau_{\rm rec}\sim\tau_{\rm mix}2^N.
 \label{eq:sm8-recurrence-time-half}
\end{equation}
The lecture uses the deliberately rough room-sized estimate
\(N\sim10^{30}\), giving the memorable scale
\(\tau_{\rm rec}/\tau_{\rm mix}\sim2^{10^{30}}\).  At ordinary atmospheric
density a room is closer to \(10^{27}\) molecules, but replacing
\(10^{30}\) by \(10^{27}\) changes none of the conclusion: the recurrence
time remains inconceivably larger than the age of the universe.

\begin{classroomqa}
\audiencequestion{In what units is the recurrence time measured?  Could a
change from seconds to hours, or to an extremely short unit, alter the
conclusion?}
\lecturerresponse{The dimensional prefactor is \(\tau_{\rm mix}\); the huge
number is the dimensionless ratio \(\tau_{\rm rec}/\tau_{\rm mix}\).
Changing seconds to hours divides by an ordinary conversion factor, which is
negligible beside \(2^{10^{30}}\).  The duration of an individual visit to
the corner does depend on molecular speeds.  The estimate concerns how rarely
the visit begins.  There is also a quantum recurrence theorem, but it is not
needed here.
\lecturetimestamp{00:17:24--00:19:18}}
\end{classroomqa}

\begin{classroomqa}
\audiencequestion{What is the objective of calculating a time that no
experiment can possibly wait for?}
\lecturerresponse{The point is conceptual.  It tells us in what sense a
closed mechanical system remains reversible.  It spends almost all its time
near equilibrium, yet on sufficiently long scales it also makes fluctuations
toward lower entropy.  The second law is therefore an overwhelmingly probable
statement about special initial data, not a microscopic equation that forbids
reverse motion.
\lecturetimestamp{00:15:12--00:16:35}}
\end{classroomqa}

\section{Reversibility and the Special Beginning}

Start knowingly in the tiny half-room region.  Almost every neighboring
trajectory immediately leaves it, because vastly more phase volume lies
outside than inside.  The gas spreads quickly and then waits an immense time
before returning.  This asymmetry between rapid relaxation and long waiting
does not violate reversible mechanics.

The time reverse of the rapid escape is not the entire waiting period.  It is
the brief, fantastically well-coordinated episode in which equilibrium-looking
molecular velocities happen to arrange themselves so that the gas
\emph{swooshes} into one half of the room.  Reverse every momentum at the end
of the ordinary escape and the gas retraces the escape in reverse.  The long
waiting time is the time required for that special incoming velocity pattern
to occur spontaneously.

\begin{classroomqa}
\audiencequestion{Does time symmetry mean that the slow recurrence into the
corner takes as long as the original escape from it?}
\lecturerresponse{No.  Compare the escape with its exact time reverse.  Both
are rapid.  What takes an enormous time is waiting in equilibrium for the
precise incoming fluctuation that initiates the reverse swoosh.  Over an
enormously long record, excursions away from equilibrium occur in both time
directions with symmetric local histories.
\lecturetimestamp{00:20:46--00:21:58}}
\end{classroomqa}

The recurrence argument consequently explains only part of the arrow of
time.  It explains why entropy ordinarily rises \emph{given} that the system
began in a low-entropy region.  It does not explain why the initial condition
was special.  In cosmology that missing premise becomes the central question:
why did the observable universe begin in such an exceptionally small region
of its gravitational phase space?

\begin{classroomqa}
\audiencequestion{Does recurrence apply to cosmology?  After the universe
spreads out, might it recur if one waits long enough?}
\lecturerresponse{That question is unavoidable, but its answer depends on
global cosmology.  A genuinely closed system with a finite equilibrium state
space invites recurrence.  An expanding spacetime, gravity, horizons, and an
infinite or nonstationary state space can change the premises.  The important
lesson here is narrower: an equilibrium recurrence picture by itself is not
a satisfactory statistical explanation of the structured universe we see.
\lecturetimestamp{00:19:59--00:20:45}}
\end{classroomqa}

\begin{classroomqa}
\audiencequestion{Does the recurrence depend on preparing the gas in a small
region at the beginning?  If the room starts at equilibrium, will the air
still eventually gather into one corner?}
\lecturerresponse{Yes.  Recurrence concerns the long-time dynamics of the
closed system, so an equilibrium-looking initial state also undergoes rare
low-entropy fluctuations.  Preparing the half-room state matters for a
different reason: it lets us condition on a known special beginning and ask
why entropy rises immediately afterwards.
\lecturetimestamp{00:22:04--00:22:20}}
\end{classroomqa}

\section{Why a Random Fluctuation Predicts Too Little}

One might try to turn recurrence into a cosmology: perhaps the universe spent
almost all of eternity in equilibrium, and our present world is one rare
fluctuation.  The proposal fails a sharp statistical test.  Condition on the
existence of an observer and ask what else that observer should expect to
see.

A fluctuation large enough to make one planet is extraordinarily rare.  A
fluctuation that makes two planets requires a larger entropy deficit and is
still rarer.  A fluctuation that makes stars, galaxies, a cosmic microwave
background, and mutually consistent records of a long past is rarer by an
enormous additional exponential.  An equilibrium-fluctuation theory
therefore predicts the smallest fluctuation sufficient for observation, not
the richly coherent world actually observed.  In its most extreme form, the
minimal observer is called a Boltzmann brain.

The logic is conditional probability.  Let \(A\) be the event that one
observer-sized fluctuation occurs and \(B\) an independent second structure.
If both have tiny probability \(p\), then
\begin{equation}
 P(B\mid A)=p,
 \qquad
 P(A\cap B)=p^2.
 \label{eq:sm8-conditional-fluctuations}
\end{equation}
Learning that the first fluctuation occurred does not make the second likely.
The lottery analogy is exact at this level: a previous winner has no improved
chance on the next independent draw.  Thus the observer may condition on
being present, but should still be astonished to find an independently
structured environment.

\begin{classroomqa}
\audiencequestion{Given that one planet or one Boltzmann observer already
exists, is the probability of another one really smaller?}
\lecturerresponse{The conditional probability for a second independent rare
fluctuation is the same tiny one-event probability, not the square of it.
The joint probability for both is the square.  This is precisely why the
first observer should still expect thermal disorder outside itself.  Finding
a spouse, another planet, and an ordered sky is evidence against the premise
that the observer was assembled as an isolated equilibrium fluctuation.
\lecturetimestamp{00:26:12--00:30:27}}
\end{classroomqa}

\subsection{Coherent records of a past}

The same argument applies to history.  Two textbooks that agree, astronomical
records that fit one expansion history, geological strata, memories, and
laboratory notebooks all constrain one another.  A single random textbook
could contain any story.  An entire network of records that agrees about the
same past occupies a much smaller phase-space region.

The classroom example uses the familiar tale that George Washington chopped
down a cherry tree.  Whether that tale is true is irrelevant.  The point is
that, in a world fluctuated into existence at the present instant, one book's
sentence gives no dynamical reason for another record to agree with it.
Coherent correlations are evidence of a shared history, not merely of one
instantaneous accident.

\begin{classroomqa}
\audiencequestion{What does \emph{historical coherence} mean here?}
\lecturerresponse{It means that many present records fit one another as
consequences of a common past.  A fluctuation could fabricate records and
memories, but fabricating a mutually consistent network requires a much
larger fluctuation than fabricating one minimal observer.  The recurrence
theory therefore predicts that most observer-like fluctuations lack the
coherence we actually find.
\lecturetimestamp{00:30:30--00:31:12}}
\end{classroomqa}

\begin{classroomqa}
\audiencequestion{Could the universe still be a closed system whose
equilibration and recurrence times are merely extremely long?}
\lecturerresponse{That possibility does not remove the statistical problem.
Over an indefinitely long history there would be many observer-like
fluctuations, and the overwhelming majority would be minimal ones without a
coherent surrounding history.  An observer selected from that ensemble
should not expect our observations.  Cosmologists therefore take the
Boltzmann-brain problem seriously as a test of cosmological measures and
late-time models, not as a claim that recurrence has already been observed.
\lecturetimestamp{00:31:13--00:32:42}}
\end{classroomqa}

\section{Life as Structure Sustained by Flow}

Life seems to pose the opposite difficulty.  Organisms maintain intricate
organization for long times; they do not look like brief accidental
fluctuations.  The resolution is that the second law applies to a complete
closed system.  A subsystem may lower or maintain its entropy provided a
larger entropy increase occurs in its environment:
\begin{equation}
 \dot S_{\rm total}
 =\dot S_{\rm subsystem}+\dot S_{\rm environment}\geq 0.
 \label{eq:sm8-open-system-second-law}
\end{equation}

The Earth is not in thermal equilibrium.  It is approximately stationary on
human timescales, but stationary does not mean static.  Energy flows through
it.  Sunlight arrives, is degraded through many processes, and is radiated
away at a lower effective temperature.  That throughput can sustain local
organization while the total entropy rises.

Water flowing through a pipe supplies the immediate mechanical image.  The
flow may contain persistent eddies and vortices.  They are organized
structures, but they are maintained by the stream.  Seal both ends, stop the
flow, and viscosity erases the eddies as the water approaches a dull
equilibrium.  The outgoing water is warmer than the incoming water, so the
larger system has paid the entropy bill.

\begin{classroomqa}
\audiencequestion{Does life decrease entropy, and why can its organization
persist rather than appear and disappear as an ordinary fluctuation?}
\lecturerresponse{Life may maintain or lower entropy locally, but only by
exporting still more entropy.  Its persistence comes from a stationary
nonequilibrium flow, chiefly the flow of solar energy through the Earth.  If
the Earth were sealed against incoming and outgoing radiation, it would
eventually approach thermal equilibrium and the organized activity would
cease.  Life is closer to an eddy in a stream than to an equilibrium
fluctuation.
\lecturetimestamp{00:32:43--00:35:42}}
\end{classroomqa}

With that, we leave the \emph{interesting things} and return, in the
lecture's deliberately dry phrase, to something dull: magnets.

\section{Cooling a Magnet: Domains and a Chosen Direction}

Statistical mechanics often calls a mathematical model a magnet even when it
is not a literal piece of iron.  The useful common feature is a collection of
local degrees of freedom that can order.  In a real magnetic material these
degrees of freedom may arise from atomic magnetic moments, electron spins,
orbital currents, or effective moments of crystal units.

At high temperature thermal disorder makes the local moments point in many
directions, giving little or no net magnetization.  If neighboring moments
have lower energy when aligned, cooling allows small aligned domains to form.
Further cooling makes the energy advantage more important and the correlated
domains grow.  In a system that supports a ferromagnetic transition, the
correlation length becomes macroscopic at a finite critical temperature.
The ordered phase then chooses an orientation even though the zero-field
energy did not prefer one orientation in advance.

\begin{classroomqa}
\audiencequestion{Should microscopic magnets prefer to anti-align because a
north pole attracts a south pole?}
\lecturerresponse{The effective interaction depends on the material.
Ordinary dipole intuition is not the whole microscopic problem; exchange,
crystal structure, and competing interactions matter.  Ferromagnetic
materials are those in which the effective low-energy tendency favors
parallel order, while other materials can favor anti-alignment.
\lecturetimestamp{00:39:56--00:40:30}}
\end{classroomqa}

\begin{classroomqa}
\audiencequestion{Must an external magnetic field be present while the
material cools in order for the moments to line up?}
\lecturerresponse{No.  The interaction can produce order without a field.
In an exactly symmetric idealization, different directions are energetically
equivalent, so the realized direction is selected by initial fluctuations,
boundaries, or an arbitrarily small stray field.  Choosing one of those
equivalent directions is spontaneous symmetry breaking; the onset of
macroscopic order is the magnetic phase transition.
\lecturetimestamp{00:40:31--00:42:33}}
\end{classroomqa}

Real materials are complicated, so we now simplify brutally.  Every local
moment may point only up or down:
\begin{equation}
 \sigma_i=+1\quad\hbox{or}\quad\sigma_i=-1.
 \label{eq:sm8-ising-spin}
\end{equation}
There is no east or west in this model.  The two allowed orientations are a
definition of the model, as heads and tails are the two outcomes of a coin.
Its value is not microscopic realism but mathematical clarity and its ability
to represent ordering phenomena shared by many systems.

\begin{classroomqa}
\audiencequestion{What are these two-state objects physically?  Should they
be imagined as particles?}
\lecturerresponse{One may picture atoms on the sites of a crystal, each with
an effective magnetic moment produced by spin or microscopic current.  For
the model, however, only the two-valued variable matters.  The restriction to
up and down is imposed to make the statistical mechanics simple.
\lecturetimestamp{00:44:22--00:45:01}}
\end{classroomqa}

\section{Independent Spins in an External Field}

The first model contains \(N\) spins that do not interact with one another.
An external field \(H\) supplies the only energy scale, and \(\mu\) measures
the strength of one spin's coupling to it.  We keep the sign convention used
in the calculation:
\begin{equation}
 \epsilon_i=\mu H\sigma_i.
 \label{eq:sm8-single-spin-energy}
\end{equation}
For \(H>0\), an up spin \(\sigma_i=+1\) has energy \(+\mu H\), while a down
spin \(\sigma_i=-1\) has energy \(-\mu H\).  Thus this convention makes the
positive field favor down spins.  It is the reverse of another common naming
convention, but no physics depends on the names if the signs are used
consistently.

Let \(n\) be the number of up spins and \(m\) the number of down spins.  Then
\begin{equation}
 n+m=N,
 \qquad
 E=(n-m)\mu H.
 \label{eq:sm8-counting-energy}
\end{equation}

\begin{figure}[t]
\centering
\includegraphics[width=0.90\textwidth]{lecture_08_figure_03.png}
\caption{The independent-spin setup: \(n\) up spins, \(m\) down spins,
\(n+m=N\), and \(E=(n-m)\mu H\).}
\lectureframe{00:49:50}
\end{figure}

\begin{classroomqa}
\audiencequestion{What are the dimensions of \(H\), \(\mu\), and \(E\)?  Is
\(H\) a number or a magnetic field?}
\lecturerresponse{\(H\) denotes the externally imposed field strength and
\(\mu\) the magnetic moment.  Their product \(\mu H\) has dimensions of
energy.  In the abstract model one often combines them into a single
energy-valued parameter, customarily also denoted by a lowercase \(h\).
Here we retain both factors to show their physical origin, and \(E\) is the
total energy.
\lecturetimestamp{00:48:21--00:49:27}}
\end{classroomqa}

\subsection{Counting configurations}

For fixed \(n\) and \(m=N-n\), choose which \(n\) of the \(N\) sites point up.
The multiplicity is
\begin{equation}
 \Omega(n,m)=\binom{N}{n}
 =\frac{N!}{n!\,m!}.
 \label{eq:sm8-spin-multiplicity}
\end{equation}
Every one of these configurations has the same energy
\eqref{eq:sm8-counting-energy}.

\begin{classroomqa}
\audiencequestion{Can \(n\) and \(m\) both be summation indices, since they
are not independent?  And where does the combinatorial factor enter?}
\lecturerresponse{They may be written as a constrained pair with
\(n+m=N\), or one may sum over \(n=0,\ldots,N\) and set \(m=N-n\).
The Boltzmann weight belongs to each individual configuration; grouping all
configurations of the same energy multiplies that weight by
\(N!/(n!\,m!)\).
\lecturetimestamp{00:52:04--00:53:03}}
\end{classroomqa}

The partition function is therefore
\begin{align}
 Z_N
 &=\sum_{n+m=N}
   \frac{N!}{n!\,m!}
   e^{-\beta\mu H(n-m)}
 \notag\\
 &=\sum_{n=0}^{N}
   \binom{N}{n}
   \left(e^{-\beta\mu H}\right)^n
   \left(e^{+\beta\mu H}\right)^{N-n}.
 \label{eq:sm8-spin-z-sum}
\end{align}
Introduce
\begin{equation}
 x=e^{-\beta\mu H},
 \qquad
 y=e^{+\beta\mu H}.
 \label{eq:sm8-xy}
\end{equation}
The binomial theorem then completes the calculation:
\begin{align}
 Z_N
 &=\sum_{n=0}^{N}\binom{N}{n}x^n y^{N-n}
   =(x+y)^N
 \notag\\
 &=\left(e^{-\beta\mu H}+e^{+\beta\mu H}\right)^N
   =\left[2\cosh(\beta\mu H)\right]^N.
 \label{eq:sm8-spin-z-closed}
\end{align}

\begin{figure}[t]
\centering
\includegraphics[width=0.90\textwidth]{lecture_08_figure_04.png}
\caption{The partition-function sum, the substitutions \(x\) and \(y\), and
the binomial reduction to a power of \(2\cosh(\beta\mu H)\).}
\lectureframe{00:56:00}
\end{figure}

There is also a shorter derivation.  Since the spins are independent, one
spin has
\begin{equation}
 z_1=e^{-\beta\mu H}+e^{+\beta\mu H}
    =2\cosh(\beta\mu H),
 \label{eq:sm8-one-spin-z}
\end{equation}
and \(Z_N=z_1^N\).  The combinatorial derivation is nevertheless worth
keeping because it shows explicitly how multiplicity competes with energy.

\section{Configuration Magnetization and Ensemble Magnetization}

For one configuration define its dimensionless spin excess per site by
\begin{equation}
 M_{\rm conf}=\frac{n-m}{N}.
 \label{eq:sm8-config-magnetization}
\end{equation}
It ranges from \(+1\) for all spins up to \(-1\) for all spins down.  The
thermodynamic magnetization used here is its Boltzmann average,
\begin{equation}
 M=\left\langle M_{\rm conf}\right\rangle
  =\left\langle\frac{n-m}{N}\right\rangle.
 \label{eq:sm8-average-magnetization}
\end{equation}
The factor \(\mu H\) does not belong in this definition.  It was briefly
written there during the calculation and then removed; \(\mu H\) converts
spin excess into energy.

Configuration by configuration,
\begin{equation}
 E=N\mu H M_{\rm conf}.
 \label{eq:sm8-energy-config-m}
\end{equation}
Linearity of averaging immediately gives
\begin{equation}
 \langle E\rangle=N\mu H M.
 \label{eq:sm8-energy-average-m}
\end{equation}

\begin{figure}[t]
\centering
\includegraphics[width=0.90\textwidth]{lecture_08_figure_05.png}
\caption{The spin row, the corrected average magnetization, and the energy
relation emphasized on the board.}
\lectureframe{01:01:22}
\end{figure}

\begin{classroomqa}
\audiencequestion{Is magnetization being assigned to a particular
configuration, or are we discussing an average magnetization?  Why does the
same relation hold for average energy?}
\lecturerresponse{Both notions appear, and they should be distinguished.
\((n-m)/N\) is the value in one configuration; \(M\) is its Boltzmann
average.  Since \(E=N\mu H M_{\rm conf}\) holds separately in every
configuration, averaging both sides with the same probabilities gives
\(\langle E\rangle=N\mu H M\).  The fixed quantities \(N\), \(\mu\), and
\(H\) pass through the average.
\lecturetimestamp{00:59:54--01:03:40}}
\end{classroomqa}

\begin{classroomqa}
\audiencequestion{How can the model align every spin north--south without
biasing up against down?}
\lecturerresponse{The absence of east and west is built into the two-state
model; it is a restriction of the allowed degrees of freedom.  The present
external-field model \emph{is} biased between up and down: with our sign
convention the field lowers the energy of down spins.  Consequently this
model has no up--down symmetry.  The zero-field Ising model introduced below
will restore that symmetry.
\lecturetimestamp{01:03:41--01:04:53}}
\end{classroomqa}

\subsection{Energy from the partition function}

The mean energy follows from the usual identity
\begin{equation}
 \langle E\rangle
 =-\frac{\partial}{\partial\beta}\log Z_N.
 \label{eq:sm8-energy-identity}
\end{equation}
Using \eqref{eq:sm8-spin-z-closed},
\begin{align}
 \log Z_N
 &=N\log 2+N\log\!\left[\cosh(\beta\mu H)\right],
 \notag\\
 \langle E\rangle
 &=-N\mu H\,
   \frac{\sinh(\beta\mu H)}{\cosh(\beta\mu H)}
 \notag\\
 &=-N\mu H\tanh(\beta\mu H).
 \label{eq:sm8-average-energy}
\end{align}
Combining this with \eqref{eq:sm8-energy-average-m} gives
\begin{equation}
 \boxed{M=-\tanh(\beta\mu H)}.
 \label{eq:sm8-magnetization-tanh}
\end{equation}

The shape of the answer contains the physics.  Recall
\begin{equation}
 \cosh x=\frac{e^x+e^{-x}}{2},
 \qquad
 \sinh x=\frac{e^x-e^{-x}}{2}.
 \label{eq:sm8-hyperbolic-definitions}
\end{equation}
Near the origin, \(\tanh x=x+O(x^3)\); for large positive \(x\), it approaches
\(1\).  Since \(\beta=1/T\),
\begin{equation}
 M\longrightarrow -1
 \quad (T\longrightarrow0),
 \qquad
 M\longrightarrow 0
 \quad (T\longrightarrow\infty).
 \label{eq:sm8-magnetization-limits}
\end{equation}
At zero temperature every spin occupies the lower-energy down state.  At
infinite temperature all configurations become equally weighted and the
field bias is washed out.  Reversing \(H\) reverses \(M\).

\begin{figure}[t]
\centering
\includegraphics[width=0.90\textwidth]{lecture_08_figure_09.png}
\caption{The hyperbolic-tangent curve approaches a constant at large
\(\beta\mu H\) and is linear near the origin.  The physical magnetization in
the lecture's convention is the negative of the plotted \(\tanh\).}
\lectureframe{01:13:14}
\end{figure}

\begin{classroomqa}
\audiencequestion{The magnetization changes continuously.  Where is the
phase-transition temperature that was mentioned earlier?}
\lecturerresponse{There is none in this model.  Independent spins in a
nonzero external field are too simple, and
\(\log Z_N=N\log[2\cosh(\beta\mu H)]\) is smooth at every finite positive
temperature.  A phase transition requires the interacting, symmetric problem
that we now approach; the two-dimensional zero-field Ising model is the first
member of this family with a finite-temperature transition.
\lecturetimestamp{01:14:47--01:15:11}}
\end{classroomqa}

\section{The Ising Model: Energy in Neighboring Pairs}

The external field chose a direction before the statistics began.  Remove
that field and let neighboring spins provide one another's effective field.
If both neighbors of a site point up, they favor up; if both point down, they
favor down; if one points each way, their effects cancel in the simplest
picture.  Energy is now associated with pairs rather than isolated sites.

For two neighboring spins, choose \(J>0\) and write
\begin{equation}
 E_{12}=-J\sigma_1\sigma_2.
 \label{eq:sm8-pair-energy}
\end{equation}
Aligned spins have \(\sigma_1\sigma_2=+1\) and energy \(-J\);
anti-aligned spins have product \(-1\) and energy \(+J\).  The minus sign is
therefore exactly what makes the interaction ferromagnetic.

\begin{figure}[t]
\centering
\includegraphics[width=0.90\textwidth]{lecture_08_figure_06.png}
\caption{The local ferromagnetic interaction
\(E_{12}=-J\sigma_1\sigma_2\).}
\lectureframe{01:19:59}
\end{figure}

On an open one-dimensional chain the Hamiltonian is
\begin{equation}
 {\cal H}=-J\sum_{i=1}^{N-1}\sigma_i\sigma_{i+1}.
 \label{eq:sm8-ising-chain-open}
\end{equation}
With periodic boundary conditions one sets \(\sigma_{N+1}=\sigma_1\) and
sums through \(N\).  Either convention counts every neighboring pair once;
their boundary difference becomes negligible for a long chain.

\begin{figure}[t]
\centering
\includegraphics[width=0.90\textwidth]{lecture_08_figure_10.png}
\caption{The local pair energy extended to the nearest-neighbor chain
Hamiltonian.}
\lectureframe{01:20:46}
\end{figure}

At infinite temperature every spin string is equally probable.  Half the
bonds are aligned on average and half anti-aligned, so
\begin{equation}
 \left\langle\sigma_i\sigma_{i+1}\right\rangle_{T=\infty}=0,
 \qquad
 \left\langle{\cal H}\right\rangle_{T=\infty}=0.
 \label{eq:sm8-infinite-t-bond}
\end{equation}
Zero energy here is not a low-energy state.  The all-aligned ground state has
energy \(-J\) per bond, so the infinite-temperature ensemble lies an
extensive energy above it.

\begin{classroomqa}
\audiencequestion{How can the average product be zero when every individual
product \(\sigma_i\sigma_{i+1}\) is either \(+1\) or \(-1\)?}
\lecturerresponse{An average need not be one of the values taken in a single
configuration.  At infinite temperature aligned and anti-aligned bonds occur
equally often, so their \(+1\) and \(-1\) contributions cancel.  The same
averaging distinction that mattered for magnetization matters here.
\lecturetimestamp{01:21:39--01:22:37}}
\end{classroomqa}

\subsection{Two ground states and one global symmetry}

At zero temperature every bond wants to be aligned.  There are consequently
two ground states,
\begin{equation}
 (+,+,+,\ldots,+)
 \qquad\hbox{and}\qquad
 (-,-,-,\ldots,-).
 \label{eq:sm8-two-ground-states}
\end{equation}
They have exactly the same energy.  The Hamiltonian makes this degeneracy
transparent: reverse every spin,
\begin{equation}
 \sigma_i\longmapsto-\sigma_i
 \qquad\hbox{for all }i,
 \label{eq:sm8-global-flip}
\end{equation}
and every bond product is unchanged,
\begin{equation}
 (-\sigma_i)(-\sigma_{i+1})
 =\sigma_i\sigma_{i+1}.
 \label{eq:sm8-global-flip-invariance}
\end{equation}
Flipping one spin is not a symmetry, because it reverses the sign of the
adjacent bond products.  Flipping all spins is a symmetry.  This is the
global \(\mathbb Z_2\) symmetry of the zero-field Ising model.

\begin{figure}[t]
\centering
\includegraphics[width=0.90\textwidth]{lecture_08_figure_11.png}
\caption{The chain Hamiltonian beside the global transformation
\(\sigma_i\mapsto-\sigma_i\), which leaves every pair product invariant.}
\lectureframe{01:26:31}
\end{figure}

Apply a tiny field to one spin at strictly zero temperature.  It breaks the
tie between the two ground states, and the ferromagnetic couplings propagate
that choice through the chain.  Remove the field after the state has been
selected and a large system can retain its orientation for a very long time.
This is the intuitive entry point to spontaneous symmetry breaking.

There is one essential statistical qualification.  For a finite system in
exactly zero field, the canonical ensemble gives equal weight to the two
ground states and hence zero average magnetization.  A broken-symmetry phase
is defined by taking the thermodynamic limit before removing the selecting
field:
\begin{equation}
 M_*=\lim_{H\to0^+}\lim_{N\to\infty}M_N(H).
 \label{eq:sm8-order-of-limits}
\end{equation}
In the one-dimensional nearest-neighbor model this limit is nonzero only at
\(T=0\); any finite temperature destroys long-range order in the infinite
chain.  In two dimensions the ordered phase survives below a nonzero
critical temperature.

\subsection{A historical correction}

The classroom account gives this history with a joke at Ising's expense.
The accurate chronology is kinder.  Wilhelm Lenz proposed the model, and his
student Ernst Ising solved the one-dimensional case in his 1924 thesis,
published in 1925.  Ising correctly found no finite-temperature
ferromagnetic transition in one dimension, although the broader significance
of higher dimensions was not yet understood.  Kramers and Wannier developed
the two-dimensional duality and identified the square-lattice critical point
in 1941; Onsager gave the celebrated exact zero-field solution in 1944.

The next calculation is now clear.  We must find the partition function of
\eqref{eq:sm8-ising-chain-open} and determine how its correlations depend on
temperature.  That calculation will show explicitly why the one-dimensional
chain has no finite-temperature phase transition and why dimensionality
matters.

\section{Where We Stand}

The recurrence problem and the spin problem are two versions of the same
statistical lesson.  In the gas, an entropy deficit makes a macrostate
exponentially rare:
\begin{equation}
 P(A)=e^{-\Delta S},
 \qquad
 \tau_{\rm rec}\sim\tau_{\rm mix}e^{\Delta S}.
 \label{eq:sm8-summary-recurrence}
\end{equation}
Reversible mechanics permits the return, while the enormous phase-volume
ratio explains why ordinary experience sees only relaxation.  The remaining
arrow-of-time problem lies in the special cosmological initial condition, not
in a failure of microscopic reversibility.

For independent spins in a field,
\begin{equation}
 Z_N=\left[2\cosh(\beta\mu H)\right]^N,
 \qquad
 M=-\tanh(\beta\mu H).
 \label{eq:sm8-summary-spins}
\end{equation}
That system is explicitly biased and changes smoothly with temperature.  The
Ising interaction
\begin{equation}
 {\cal H}=-J\sum_i\sigma_i\sigma_{i+1}
 \label{eq:sm8-summary-ising}
\end{equation}
removes the external bias, restores a global up--down symmetry, and presents
the real ordering question: can a symmetric many-body system choose one of
two ordered states at nonzero temperature?
